\scanpage{48}\mspage{XIX}
tou\deleted{\unsure{s}}\unsure{tes|ts} les \deleted{\unsure{so|sa}}\unsure{.} arêtes de $\Lambda^{\unsure{1'|11}}$ sauf une (savoir, $[z,s]$) appartiennent à $\deleted{\unsure{A|R}}\,A(\Gamma)$ et le cardinal de vérifie $A(\Lambda^{1'})\cap A(\Gamma)<A(\Gamma)\cap A(\Lambda^1)$.

Comme il \unsure{n'ya|n'y a} qu'un nombre fini d'arêtes de $\Gamma$ il est clair qu'on \deleted{\unsure{bien}} a :

a) Il existe $n$ tel que $\deleted{\unsure{A|R}}\,\Lambda^{1(n)}$ (Notation definie par \unsure{récurrence|récurence} de \unsure{foçon|façon} \unsure{evidente|évidente}) contienne un \inserted{bon} sommet $t$ tel que les deux arêtes qui contiennent $t$ soient des arêtes $\Gamma$, et par conséquent $t$ est aussi un bon sommet de $\Gamma$ compte tenu du lemme $\overline{\underline{\mathrm{XIX}}}$, appliqué $n$-fois, \inserted{ou bien,}

b) il existe $n$ tel que $\mathrm{card}\,(A({\unsure{\Lambda|A}}^{1(n)})\cap A(\Gamma))=2$ donc $\mathrm{card}\,A^{1(n)}=3$ et il est clair que le \inserted{seul} sommet de $\Lambda^{1(n)}$ qui appartient aux deux arêtes de $\Lambda^{1(n)}$ contenues dans $A(\Gamma)$ est un bon sommet de $\Gamma$, compte tenu encore une fois du lemme $\overline{\underline{\mathrm{XIX}}}$\inserted{appliqué $n$ fois\unsure{,|.} Le même argument avec $\Lambda^2$ prouve le reste,}, \unsure{c.q.f.d.|c.q.s.d.}

\uline{Lemme $\overline{\underline{\mathrm{XXI}}}$} : Si $\Gamma$ est un polygone non-convexe il existe un bon sommet de $\Gamma$ tel que le \inserted{\unsure{interieur|intérieur} du} triangle corespondant est contenu dans $\mathrm{ext}\,\Gamma$.

\uline{Demonstration} : le lemme $\overline{\underline{\mathrm{VIII}}}$ montre \unsure{qu il|qu'il} \unsure{exìste|existe} un sommet $s$ de $\Gamma$ et un sommet $s'$ visible de $s$ par l'\unsure{exterieur|extérieur} de $\Gamma$. On conserve ici les notations du \unsure{numèro|numéro|numero} $3$. On considère $\Lambda^2$ (defini comme dans le numéro $3$) il est clair \margin{\unsure{faux}, \unsure{si on} \unsure{ne choisit} \unsure{pas bien} \unsure{ses|s.} \unsure{notations|solutions} \unsure{--|…}} % slanted
que $\mathrm{int}\,\Lambda^2\subset\mathrm{ext}\,\Gamma$ car $V=\mathrm{ext}\,\Lambda^1\cup\mathrm{int}\,\Lambda^2\cup{]s,s'[}$, et $V$ est non \deleted{\unsure{co|lo}}-relativement compact. \deleted{L'argument du}\inserted{Le} lemme \unsure{précedent|précédent} appliqué a $\Lambda^2$ montre que l'on peut trouver un \deleted{\unsure{trian|triang}} bon sommet $s$ de $\Lambda^2$ tel que : le triangle \unsure{corespondent|corespondant} soit contenu dans $\mathrm{ext}\,\Gamma$ et les deux arêtes de $\Lambda^2$ auxquelles \margin{\uline{\unsure{douteux|donc ?}}} % slanted
$s$ appartient soient des arêtes de $\Gamma$, ce qui permet de prouver le lemme.

\deleted{\unsure{Le|Les}}\unsure{Le|Lee} lemme $\overline{\underline{\mathrm{XX}}}$ le lemme $\overline{\underline{\mathrm{XXI}}}$ joints au fait que dans tout polygone convexe il 'y a au moins trois bons sommets, \unsure{montent|montrent} le théorème $2$, c.q.f.d.

\begin{flushright}
le 15 \unsure{juillet|jullet} 1976, Montpellier\\
\unsure{Mauré|Mouré|Maury|Maurice}.
\end{flushright}
% archive sheet number: 47
